\documentclass[10pt,a4paper]{amsart}
\usepackage{amsmath,amssymb,amsthm,graphicx}
% For commutative diagrams we use Xy-pic
\usepackage[all]{xy}

% We use 2cell for 2-commutative diagrams.
\xyoption{2cell}
% Theorem environments.
%
\theoremstyle{plain}
\newtheorem{theorem}[subsection]{Theorem}
\newtheorem{proposition}[subsection]{Proposition}
\newtheorem{lemma}[subsection]{Lemma}

\theoremstyle{definition}
\newtheorem{definition}[subsection]{Definition}
\newtheorem{example}[subsection]{Example}
\newtheorem{exercise}[subsection]{Exercise}
\newtheorem{situation}[subsection]{Situation}

\theoremstyle{remark}
\newtheorem{remark}[subsection]{Remark}
\newtheorem{remarks}[subsection]{Remarks}

\numberwithin{equation}{subsection}

% Macros
%
\def\lim{\mathop{\mathrm{lim}}\nolimits}
\def\colim{\mathop{\mathrm{colim}}\nolimits}
\def\Spec{\mathop{\mathrm{Spec}}}
\def\Hom{\mathop{\mathrm{Hom}}\nolimits}
\def\Ext{\mathop{\mathrm{Ext}}\nolimits}
\def\SheafHom{\mathop{\mathcal{H}\!\mathit{om}}\nolimits}
\def\SheafExt{\mathop{\mathcal{E}\!\mathit{xt}}\nolimits}
\def\Sch{\mathit{Sch}}
\def\Mor{\mathop{\mathrm{Mor}}\nolimits}
\def\Ob{\mathop{\mathrm{Ob}}\nolimits}
\def\Sh{\mathop{\mathit{Sh}}\nolimits}
\def\NL{\mathop{N\!L}\nolimits}
\def\CH{\mathop{\mathrm{CH}}\nolimits}
\def\proetale{{pro\text{-}\acute{e}tale}}
\def\etale{{\acute{e}tale}}
\def\QCoh{\mathit{QCoh}}
\def\Ker{\mathop{\mathrm{Ker}}}
\def\Im{\mathop{\mathrm{Im}}}
\def\Coker{\mathop{\mathrm{Coker}}}
\def\Coim{\mathop{\mathrm{Coim}}}

% Boxtimes
%
\DeclareMathSymbol{\boxtimes}{\mathbin}{AMSa}{"02}

%
% Macros for moduli stacks/spaces
%
\def\QCohstack{\mathcal{QC}\!\mathit{oh}}
\def\Cohstack{\mathcal{C}\!\mathit{oh}}
\def\Spacesstack{\mathcal{S}\!\mathit{paces}}
\def\Quotfunctor{\mathrm{Quot}}
\def\Hilbfunctor{\mathrm{Hilb}}
\def\Curvesstack{\mathcal{C}\!\mathit{urves}}
\def\Polarizedstack{\mathcal{P}\!\mathit{olarized}}
\def\Complexesstack{\mathcal{C}\!\mathit{omplexes}}
% \Pic is the operator that assigns to X its picard group, usage \Pic(X)
% \Picardstack_{X/B} denotes the Picard stack of X over B
% \Picardfunctor_{X/B} denotes the Picard functor of X over B
\def\Pic{\mathop{\mathrm{Pic}}\nolimits}
\def\Picardstack{\mathcal{P}\!\mathit{ic}}
\def\Picardfunctor{\mathrm{Pic}}
\def\Deformationcategory{\mathcal{D}\!\mathit{ef}}
\usepackage{cleveref}
\crefname{originallemma}{original lemma}{original lemmas}
\Crefname{originallemma}{Original lemma}{Original lemmas}
\makeatletter
\newlabel{derived-lemma-adjoint-preserve-K-injectives}{{08BJ}{source}{}{sourcepointer-1}{}}
\newlabel{derived-lemma-adjoint-preserve-K-injectives@cref}{{[originallemma][0][]08BJ}{[1][0][]source}{}{}{}}
\makeatother
\crefname{originallemma}{original lemma}{original lemmas}
\Crefname{originallemma}{Original lemma}{Original lemmas}
\makeatletter
\newlabel{homology-lemma-adjoint-preserve-injectives}{{015Z}{source}{}{sourcepointer-2}{}}
\newlabel{homology-lemma-adjoint-preserve-injectives@cref}{{[originallemma][0][]015Z}{[1][0][]source}{}{}{}}
\makeatother
\crefname{originallemma}{original lemma}{original lemmas}
\Crefname{originallemma}{Original lemma}{Original lemmas}
\makeatletter
\newlabel{sites-cohomology-lemma-cohomology-of-open}{{03F3}{source}{}{sourcepointer-3}{}}
\newlabel{sites-cohomology-lemma-cohomology-of-open@cref}{{[originallemma][0][]03F3}{[1][0][]source}{}{}{}}
\makeatother
\crefname{originallemma}{original lemma}{original lemmas}
\Crefname{originallemma}{Original lemma}{Original lemmas}
\makeatletter
\newlabel{sites-cohomology-lemma-restrict-K-injective-to-open}{{08FI}{source}{}{sourcepointer-4}{}}
\newlabel{sites-cohomology-lemma-restrict-K-injective-to-open@cref}{{[originallemma][0][]08FI}{[1][0][]source}{}{}{}}
\makeatother
\crefname{originallemma}{original lemma}{original lemmas}
\Crefname{originallemma}{Original lemma}{Original lemmas}
\makeatletter
\newlabel{sites-modules-lemma-exactness-lower-shriek}{{04BH}{source}{}{sourcepointer-5}{}}
\newlabel{sites-modules-lemma-exactness-lower-shriek@cref}{{[originallemma][0][]04BH}{[1][0][]source}{}{}{}}
\makeatother
\crefname{originallemma}{original lemma}{original lemmas}
\Crefname{originallemma}{Original lemma}{Original lemmas}
\makeatletter
\newlabel{sites-modules-lemma-g-shriek-adjoint}{{04BG}{source}{}{sourcepointer-6}{}}
\newlabel{sites-modules-lemma-g-shriek-adjoint@cref}{{[originallemma][0][]04BG}{[1][0][]source}{}{}{}}
\makeatother
\crefname{originallemma}{original lemma}{original lemmas}
\Crefname{originallemma}{Original lemma}{Original lemmas}
\makeatletter
\newlabel{topologies-lemma-morphism-big-small-etale}{{021I}{source}{}{sourcepointer-7}{}}
\newlabel{topologies-lemma-morphism-big-small-etale@cref}{{[originallemma][0][]021I}{[1][0][]source}{}{}{}}
\makeatother
\crefname{originallemma}{original lemma}{original lemmas}
\Crefname{originallemma}{Original lemma}{Original lemmas}
\makeatletter
\newlabel{topologies-lemma-put-in-T-etale}{{021F}{source}{}{sourcepointer-8}{}}
\newlabel{topologies-lemma-put-in-T-etale@cref}{{[originallemma][0][]021F}{[1][0][]source}{}{}{}}
\makeatother
\title{Injectivity and K-injectivity survive restriction of abelian sheaves from the big etale site to a small one, but fail for modules}
\author{The Stacks Project Authors}
\date{stacks-project:etale-cohomology.tex:local-snapshot-20261011; GNU FDL 1.2 or later}
\begin{document}
\maketitle
\noindent\textit{Editorial note.} All delivered mathematical text is the author's own, copied from the pinned Stacks Project chapter file etale-cohomology.tex; no statement or proof step was rewritten, reordered or removed.\par\medskip
\noindent\textit{Editorial note.} The counted proof is the author's complete argument: the proof is written out in full: the adjointness arguments for the small and big site parts (the latter with two independent proofs) and the explicit multiplication-by-2 counterexample showing that the module statements fail.\par\medskip
\noindent\textit{Editorial note.} The label lemma-compare-injectives names the counted statement, and the proof environment immediately following it in the source is delivered verbatim as its proof.\par\medskip
\noindent\textit{Editorial note.} The counted statement is carried by the author's own lemma environment in the chapter source.\par\medskip
\noindent\textit{Editorial note.} The statement and proof are copied from the pinned chapter file; the eight cited inputs are kept as external pointers with their canonical Stacks tags.\par\medskip
\noindent\textit{Editorial note.} The preamble is the chapter's own native preamble excerpt already used for this source: the Xy-pic setup (native preamble lines 14-18) and the native theorem declarations with the mathematical macros (native preamble lines 154-217).\par\medskip
\noindent\textit{Allowed standard prerequisites (editorial).}\begin{itemize}
\item The restriction functors i\_f\textasciicircum{}\{-1\} and pi\_\{S,*\} = i\_S\textasciicircum{}\{-1\} with their exact left adjoints. Same chapter, the functors i\_f : T\_etale -> (Sch/S)\_etale and their lower shriek adjoints i\_\{f,!\}, exact by the modules-on-sites exactness result. Supplies the right-adjoint presentation that makes the two functors preserve injectives and K-injectives.
\item Right adjoints of exact functors preserve injectives and K-injectives. Homology and Derived Categories, the two lemmas stating preservation of injectives and K-injectives for a right adjoint of an exact functor. Supplies the formal statement used for parts (1)(b) and (2)(b) and again for the localisation step in (1)(a) and (2)(a).
\item Localisation morphisms preserve injectives and K-injectives. Cohomology on Sites, the lemmas on restriction to an open and K-injective complexes on localisations. Supplies the second proof of parts (1)(a) and (2)(a), through the factorisation i\_f = i\_T composed with f\_big.
\end{itemize}\medskip
\ifdefined\hypertarget\hypertarget{sourcepointer-1}{}\fi
\noindent\textit{Source reference.} Original lemma 08BJ: sources/stacks-chow/tags.txt, source lines 10569--10569. Prerequisite pointer invoked at proof line(s) 21128; it names a standard fact of a sibling chapter that the counted argument uses, and every step of the counted proof built on it is retained verbatim in the delivered body.\par\medskip
\ifdefined\hypertarget\hypertarget{sourcepointer-2}{}\fi
\noindent\textit{Source reference.} Original lemma 015Z: sources/stacks-chow/tags.txt, source lines 1526--1526. Prerequisite pointer invoked at proof line(s) 21127; it names a standard fact of a sibling chapter that the counted argument uses, and every step of the counted proof built on it is retained verbatim in the delivered body.\par\medskip
\ifdefined\hypertarget\hypertarget{sourcepointer-3}{}\fi
\noindent\textit{Source reference.} Original lemma 03F3: sources/stacks-chow/tags.txt, source lines 4478--4478. Prerequisite pointer invoked at proof line(s) 21143; it names a standard fact of a sibling chapter that the counted argument uses, and every step of the counted proof built on it is retained verbatim in the delivered body.\par\medskip
\ifdefined\hypertarget\hypertarget{sourcepointer-4}{}\fi
\noindent\textit{Source reference.} Original lemma 08FI: sources/stacks-chow/tags.txt, source lines 10711--10711. Prerequisite pointer invoked at proof line(s) 21144; it names a standard fact of a sibling chapter that the counted argument uses, and every step of the counted proof built on it is retained verbatim in the delivered body.\par\medskip
\ifdefined\hypertarget\hypertarget{sourcepointer-5}{}\fi
\noindent\textit{Source reference.} Original lemma 04BH: sources/stacks-chow/tags.txt, source lines 5652--5652. Prerequisite pointer invoked at proof line(s) 21138; it names a standard fact of a sibling chapter that the counted argument uses, and every step of the counted proof built on it is retained verbatim in the delivered body.\par\medskip
\ifdefined\hypertarget\hypertarget{sourcepointer-6}{}\fi
\noindent\textit{Source reference.} Original lemma 04BG: sources/stacks-chow/tags.txt, source lines 5651--5651. Prerequisite pointer invoked at proof line(s) 21136; it names a standard fact of a sibling chapter that the counted argument uses, and every step of the counted proof built on it is retained verbatim in the delivered body.\par\medskip
\ifdefined\hypertarget\hypertarget{sourcepointer-7}{}\fi
\noindent\textit{Source reference.} Original lemma 021I: sources/stacks-chow/tags.txt, source lines 2671--2671. Prerequisite pointer invoked at proof line(s) 21140; it names a standard fact of a sibling chapter that the counted argument uses, and every step of the counted proof built on it is retained verbatim in the delivered body.\par\medskip
\ifdefined\hypertarget\hypertarget{sourcepointer-8}{}\fi
\noindent\textit{Source reference.} Original lemma 021F: sources/stacks-chow/tags.txt, source lines 2668--2668. Prerequisite pointer invoked at proof line(s) 21134; it names a standard fact of a sibling chapter that the counted argument uses, and every step of the counted proof built on it is retained verbatim in the delivered body.\par\medskip
\begin{lemma}
\label{lemma-compare-injectives}
Let $S$ be a scheme. Let $T$ be an object of $(\Sch/S)_\etale$.
\begin{enumerate}
\item If $\mathcal{I}$ is injective in $\textit{Ab}((\Sch/S)_\etale)$, then
\begin{enumerate}
\item $i_f^{-1}\mathcal{I}$ is injective in $\textit{Ab}(T_\etale)$,
\item $\mathcal{I}|_{S_\etale}$ is injective in $\textit{Ab}(S_\etale)$,
\end{enumerate}
\item If $\mathcal{I}^\bullet$ is a K-injective complex
in $\textit{Ab}((\Sch/S)_\etale)$, then
\begin{enumerate}
\item $i_f^{-1}\mathcal{I}^\bullet$ is a K-injective complex in
$\textit{Ab}(T_\etale)$,
\item $\mathcal{I}^\bullet|_{S_\etale}$ is a K-injective complex in
$\textit{Ab}(S_\etale)$,
\end{enumerate}
\end{enumerate}
The corresponding statements for modules do not hold.
\end{lemma}

\begin{proof}
Parts (1)(b) and (2)(b)
follow formally from the fact that the restriction functor
$\pi_{S, *} = i_S^{-1}$ is a right adjoint of the exact functor
$\pi_S^{-1}$, see
Homology, Lemma \ref{homology-lemma-adjoint-preserve-injectives} and
Derived Categories, Lemma \ref{derived-lemma-adjoint-preserve-K-injectives}.

\medskip\noindent
Parts (1)(a) and (2)(a) can be seen in two ways. First proof: We can use
that $i_f^{-1}$ is a right adjoint of the exact functor $i_{f, !}$.
This functor is constructed in
Topologies, Lemma \ref{topologies-lemma-put-in-T-etale}
for sheaves of sets and for abelian sheaves in
Modules on Sites, Lemma \ref{sites-modules-lemma-g-shriek-adjoint}.
It is shown in Modules on Sites, Lemma
\ref{sites-modules-lemma-exactness-lower-shriek} that it is exact.
Second proof. We can use that $i_f = i_T \circ f_{big}$ as is shown
in Topologies, Lemma \ref{topologies-lemma-morphism-big-small-etale}.
Since $f_{big}$ is a localization, we see that pullback by it
preserves injectives and K-injectives, see
Cohomology on Sites, Lemmas \ref{sites-cohomology-lemma-cohomology-of-open} and
\ref{sites-cohomology-lemma-restrict-K-injective-to-open}.
Then we apply the already proved parts (1)(b) and (2)(b)
to the functor $i_T^{-1}$ to conclude.

\medskip\noindent
Let $S = \Spec(\mathbf{Z})$ and consider the map
$2 : \mathcal{O}_S \to \mathcal{O}_S$. This is an injective map
of $\mathcal{O}_S$-modules on $S_\etale$. However, the pullback
$\pi_S^*(2) : \mathcal{O} \to \mathcal{O}$ is not injective
as we see by evaluating on $\Spec(\mathbf{F}_2)$. Now choose
an injection $\alpha : \mathcal{O} \to \mathcal{I}$ into an injective
$\mathcal{O}$-module $\mathcal{I}$ on $(\Sch/S)_\etale$.
Then consider the diagram
$$
\xymatrix{
\mathcal{O}_S \ar[d]_2 \ar[rr]_{\alpha|_{S_\etale}} & &
\mathcal{I}|_{S_\etale} \\
\mathcal{O}_S \ar@{..>}[rru]
}
$$
Then the dotted arrow cannot exist in the category of $\mathcal{O}_S$-modules
because it would mean
(by adjunction) that the injective map $\alpha$ factors through
the noninjective map $\pi_S^*(2)$ which cannot be the case.
Thus $\mathcal{I}|_{S_\etale}$ is not an injective $\mathcal{O}_S$-module.
\end{proof}

\end{document}
